% GENERATED FILE. Edit canonical lessons, then run npm run generate:books.
% canonical-source-hash: fnv1a64:f346470791c2ab6e
% canonical-lessons: GU-C170-mithai, GU-C170-ladu, GU-C170-makhan, GU-C170-calvu, GU-C170-cotvu, GU-R170-review-165-170

\chapter{Laddoos and butter}
\label{ch:gu-laddoos-and-butter}
\hlchaptermodality{170}

\begin{chapteropening}
\textbf{By the end of this chapter:} \emph{I can talk about sweets, butter and sifting flour.}

Complete the last lesson of chapter 170: I can talk about sweets, butter and sifting flour.
\end{chapteropening}

\section[\texorpdfstring{\gu{મીઠાઈ} (mīṭhāī)}{મીઠાઈ (mīṭhāī)}]{\gu{મીઠાઈ} (mīṭhāī) — sweets, sweetmeat}

\label{lesson:GU-C170-mithai}



\begin{quote}
Before the new one: say the Gujarati for wok, then the Gujarati for lid.
\end{quote}

\subsection*{You'll want to know: \gu{મીઠાઈ}}
\textbf{\gu{મીઠાઈ}} — \emph{mīṭhāī} — \textquotedblleft{}sweets, sweetmeat\textquotedblright{}.

\begin{quote}
\textbf{\gu{દિવાળીમાં} \gu{અમે} \gu{મીઠાઈ} \gu{ખાઈએ} \gu{છીએ}.} — \emph{divāḷīmā̃ ame mīṭhāī khāīe chīe} — \textquotedblleft{}We eat sweets at Diwali.\textquotedblright{}
\end{quote}

Built on \textbf{\gu{મીઠું}}, \emph{mīṭhũ}, sweet.

\begin{cousinweb}
\textbf{\gu{મીઠાઈ}} (\emph{mīṭhāī}) is built on the inherited \emph{mīṭhũ}, \textquotedblleft{}sweet\textquotedblright{}, with the noun ending \emph{-āī}; EtymDB derives it from that word. CDIAL traces \emph{mīṭhũ} to Sanskrit \emph{miṣṭa}, \textquotedblleft{}dainty, sweet\textquotedblright{} (CDIAL 10299).
\end{cousinweb}

\begin{grammarlens}[title={What you've built}]
One more word for everyday life at A2.
\end{grammarlens}

\subsection*{Your turn}
\begin{itemize}
\raggedright
  \item \emph{Say it:} \emph{mīṭhāī}
  \item \emph{Say it:} \emph{mīṭhāī}, once more
  \item \emph{Say it:} \emph{divāḷīmā̃ ame mīṭhāī khāīe chīe} — We eat sweets at Diwali.
  \item \emph{Recall:} say the Gujarati for wok, then the Gujarati for lid, then say \emph{mīṭhāī} again
\end{itemize}

\begin{tcolorbox}[breakable,colback=teal!4,colframe=teal!35!black,title={Before you move on}]
What does \gu{મીઠાઈ} mean? (\textquotedblleft{}Sweets, sweetmeat\textquotedblright{}.) Say it once more.
\end{tcolorbox}
\section[\texorpdfstring{\gu{લાડુ} (lāḍu)}{લાડુ (lāḍu)}]{\gu{લાડુ} (lāḍu) — laddoo, sweet ball}

\label{lesson:GU-C170-ladu}



\begin{quote}
Before the new one: say the Gujarati for lid, then the Gujarati for sweets.
\end{quote}

\subsection*{You'll want to know: \gu{લાડુ}}
\textbf{\gu{લાડુ}} — \emph{lāḍu} — \textquotedblleft{}laddoo, sweet ball\textquotedblright{}.

\begin{quote}
\textbf{\gu{દાદી} \gu{લાડુ} \gu{બનાવે} \gu{છે}.} — \emph{dādī lāḍu banāve che} — \textquotedblleft{}Grandmother makes laddoos.\textquotedblright{}
\end{quote}

\begin{cousinweb}
\textbf{\gu{લાડુ}} (\emph{lāḍu}) is a tadbhava. It comes from Sanskrit \emph{laḍḍu}, \textquotedblleft{}a kind of sweetmeat\textquotedblright{}, and CDIAL glosses the Gujarati word \textquotedblleft{}a sweetmeat of sweetened meal in a ball\textquotedblright{} (CDIAL 10926).
\end{cousinweb}

\begin{grammarlens}[title={What you've built}]
One more word for everyday life at A2.
\end{grammarlens}

\subsection*{Your turn}
\begin{itemize}
\raggedright
  \item \emph{Say it:} \emph{lāḍu}
  \item \emph{Say it:} \emph{lāḍu}, once more
  \item \emph{Say it:} \emph{dādī lāḍu banāve che} — Grandmother makes laddoos.
  \item \emph{Recall:} say the Gujarati for lid, then the Gujarati for sweets, then say \emph{lāḍu} again
\end{itemize}

\begin{tcolorbox}[breakable,colback=teal!4,colframe=teal!35!black,title={Before you move on}]
What does \gu{લાડુ} mean? (\textquotedblleft{}Laddoo, sweet ball\textquotedblright{}.) Say it once more.
\end{tcolorbox}
\section[\texorpdfstring{\gu{માખણ} (mākhaṇ)}{માખણ (mākhaṇ)}]{\gu{માખણ} (mākhaṇ) — butter}

\label{lesson:GU-C170-makhan}



\begin{quote}
Before the new one: say the Gujarati for sweets, then the Gujarati for laddoo.
\end{quote}

\subsection*{You'll want to know: \gu{માખણ}}
\textbf{\gu{માખણ}} — \emph{mākhaṇ} — \textquotedblleft{}butter\textquotedblright{}.

\begin{quote}
\textbf{\gu{બાળકને} \gu{માખણ} \gu{ગમે} \gu{છે}.} — \emph{bāḷakne mākhaṇ game che} — \textquotedblleft{}The child likes butter.\textquotedblright{}
\end{quote}

\begin{cousinweb}
\textbf{\gu{માખણ}} (\emph{mākhaṇ}) is a tadbhava. It comes from Sanskrit \emph{mrakṣaṇa}, \textquotedblleft{}rubbing in; ointment, oil\textquotedblright{}, and CDIAL glosses the Gujarati word \textquotedblleft{}butter, ointment\textquotedblright{} (CDIAL 10378).
\end{cousinweb}

\begin{grammarlens}[title={What you've built}]
One more word for everyday life at A2.
\end{grammarlens}

\subsection*{Your turn}
\begin{itemize}
\raggedright
  \item \emph{Say it:} \emph{mākhaṇ}
  \item \emph{Say it:} \emph{mākhaṇ}, once more
  \item \emph{Say it:} \emph{bāḷakne mākhaṇ game che} — The child likes butter.
  \item \emph{Recall:} say the Gujarati for sweets, then the Gujarati for laddoo, then say \emph{mākhaṇ} again
\end{itemize}

\begin{tcolorbox}[breakable,colback=teal!4,colframe=teal!35!black,title={Before you move on}]
What does \gu{માખણ} mean? (\textquotedblleft{}Butter\textquotedblright{}.) Say it once more.
\end{tcolorbox}
\section[\texorpdfstring{\gu{ચાળવું} (cāḷvũ)}{ચાળવું (cāḷvũ)}]{\gu{ચાળવું} (cāḷvũ) — to sift}

\label{lesson:GU-C170-calvu}



\begin{quote}
Before the new one: say the Gujarati for laddoo, then the Gujarati for butter.
\end{quote}

\subsection*{You'll want to know: \gu{ચાળવું}}
\textbf{\gu{ચાળવું}} — \emph{cāḷvũ} — \textquotedblleft{}to sift\textquotedblright{}.

\begin{quote}
\textbf{\gu{તે} \gu{લોટ} \gu{ચાળે} \gu{છે}.} — \emph{te loṭ cāḷe che} — \textquotedblleft{}She sifts the flour.\textquotedblright{}
\end{quote}

\begin{cousinweb}
\textbf{\gu{ચાળવું}} (\emph{cāḷvũ}) is a tadbhava. It comes from Sanskrit \emph{cālayati}, \textquotedblleft{}shakes, jogs\textquotedblright{}, and CDIAL glosses the Gujarati verb \textquotedblleft{}to sift\textquotedblright{} (CDIAL 4772).
\end{cousinweb}

\begin{grammarlens}[title={What you've built}]
One more word for everyday life at A2.
\end{grammarlens}

\subsection*{Your turn}
\begin{itemize}
\raggedright
  \item \emph{Say it:} \emph{cāḷvũ}
  \item \emph{Say it:} \emph{cāḷvũ}, once more
  \item \emph{Say it:} \emph{te loṭ cāḷe che} — She sifts the flour.
  \item \emph{Recall:} say the Gujarati for laddoo, then the Gujarati for butter, then say \emph{cāḷvũ} again
\end{itemize}

\begin{tcolorbox}[breakable,colback=teal!4,colframe=teal!35!black,title={Before you move on}]
What does \gu{ચાળવું} mean? (\textquotedblleft{}To sift\textquotedblright{}.) Say it once more.
\end{tcolorbox}
\section[\texorpdfstring{\gu{ચોંટવું} (cõṭvũ)}{ચોંટવું (cõṭvũ)}]{\gu{ચોંટવું} (cõṭvũ) — to stick}

\label{lesson:GU-C170-cotvu}



\begin{quote}
Before the new one: say the Gujarati for butter, then the Gujarati for to sift.
\end{quote}

\subsection*{You'll want to know: \gu{ચોંટવું}}
\textbf{\gu{ચોંટવું}} — \emph{cõṭvũ} — \textquotedblleft{}to stick\textquotedblright{}.

\begin{quote}
\textbf{\gu{ભાત} \gu{તપેલીમાં} \gu{ચોંટે} \gu{છે}.} — \emph{bhāt tapelīmā̃ cõṭe che} — \textquotedblleft{}The rice sticks to the pot.\textquotedblright{}
\end{quote}

\begin{cousinweb}
The origin of \textbf{\gu{ચોંટવું}} (\emph{cõṭvũ}) is not settled. CDIAL lists no Gujarati form of it with the sense of sticking, and EtymDB has no entry for it.
\end{cousinweb}

\begin{grammarlens}[title={What you've built}]
One more word for everyday life at A2.
\end{grammarlens}

\subsection*{Your turn}
\begin{itemize}
\raggedright
  \item \emph{Say it:} \emph{cõṭvũ}
  \item \emph{Say it:} \emph{cõṭvũ}, once more
  \item \emph{Say it:} \emph{bhāt tapelīmā̃ cõṭe che} — The rice sticks to the pot.
  \item \emph{Recall:} say the Gujarati for butter, then the Gujarati for to sift, then say \emph{cõṭvũ} again
\end{itemize}

\begin{tcolorbox}[breakable,colback=teal!4,colframe=teal!35!black,title={Before you move on}]
What does \gu{ચોંટવું} mean? (\textquotedblleft{}To stick\textquotedblright{}.) Say it once more.
\end{tcolorbox}
\section[(dialogue)]{Words from chapters 165–170 — a review}

\label{lesson:GU-R170-review-165-170}



\begin{quote}
Say the Gujarati for to add and to pour.
\end{quote}

\subsection*{Your turn}
Say these aloud:

\begin{itemize}
\raggedright
  \item to add, to pour, to stir, to mix, to boil
  \item to chop, to peel, onion, fresh coriander, to grind
  \item okra, pea, egg, to steam, to fry
  \item to knead, rolling pin, to roll out, griddle, to roast
  \item bhakhri, khichdi, papad, wok, lid
  \item sweets, laddoo, butter, to sift, to stick
\end{itemize}

\begin{tcolorbox}[breakable,colback=teal!4,colframe=teal!35!black,title={Before you move on}]
To add? (\textbf{\gu{ઉમેરવું}}.) To pour? (\textbf{\gu{રેડવું}}.) To stir? (\textbf{\gu{હલાવવું}}.) To mix? (\textbf{\gu{ભેળવવું}}.) To boil? (\textbf{\gu{ઉકાળવું}}.) To chop? (\textbf{\gu{સમારવું}}.) To peel? (\textbf{\gu{છોલવું}}.) Onion? (\textbf{\gu{ડુંગળી}}.) Fresh coriander? (\textbf{\gu{કોથમીર}}.) To grind? (\textbf{\gu{વાટવું}}.) Okra? (\textbf{\gu{ભીંડો}}.) Pea? (\textbf{\gu{વટાણો}}.) Egg? (\textbf{\gu{ઈંડું}}.) To steam? (\textbf{\gu{બાફવું}}.) To fry? (\textbf{\gu{તળવું}}.) To knead? (\textbf{\gu{મસળવું}}.) Rolling pin? (\textbf{\gu{વેલણ}}.) To roll out? (\textbf{\gu{વણવું}}.) Griddle? (\textbf{\gu{તવો}}.) To roast? (\textbf{\gu{શેકવું}}.) Bhakhri? (\textbf{\gu{ભાખરી}}.) Khichdi? (\textbf{\gu{ખીચડી}}.) Papad? (\textbf{\gu{પાપડ}}.) Wok? (\textbf{\gu{કડાઈ}}.) Lid? (\textbf{\gu{ઢાંકણું}}.) Sweets? (\textbf{\gu{મીઠાઈ}}.) Laddoo? (\textbf{\gu{લાડુ}}.) Butter? (\textbf{\gu{માખણ}}.) To sift? (\textbf{\gu{ચાળવું}}.) To stick? (\textbf{\gu{ચોંટવું}}.)
\end{tcolorbox}
